\documentclass[10pt,a4paper]{amsart}
\usepackage[margin=19mm]{geometry}
\usepackage{amsmath,amssymb,mathtools}
\usepackage[scr=rsfs,cal=euler]{mathalfa}
\usepackage{xfrac}
\usepackage[unicode,hidelinks,bookmarks=false]{hyperref}
\newcommand{\ie}{i.e.\ }
\newcommand{\cf}{cf.\ }
\newcommand{\resp}{respectively\ }
\newcommand{\Subsection}{\subsection}
\providecommand{\nobreakdash}{\nobreak\hbox{-}}
\setlength{\emergencystretch}{2em}
\multlinegap0pt
\usepackage{xcolor}
\usepackage{amssymb}
\usepackage[shortlabels]{enumitem}
\newtheorem{lemma}{Lemma}
\newtheorem{proposition}{Proposition}
\newcommand{\E}{\mathbb{E}}
\newcommand{\R}{\mathbb{R}}
\DeclareMathOperator*{\argmin}{arg\,min}
\def\sU{{\mathsf U}}
\def\sX{{\mathsf X}}
\def\sZ{{\mathsf Z}}
\title{One if by Land, Two if by Sea, Three if by Four Seas, and More to Come -- Values of Perception, Prediction, Communication, and Common Sense in Decision Making}
\author{Aolin Xu}
\date{Source: arXiv:2601.06077v2 (CC BY 4.0)}
\begin{document}
\maketitle

\noindent\emph{Source and licence.} One if by Land, Two if by Sea, Three if by Four Seas, and More to Come -- Values of Perception, Prediction, Communication, and Common Sense in Decision Making. Version arXiv:2601.06077v2 (CC BY 4.0), distributed by arXiv under the Creative Commons Attribution 4.0 International licence, http://creativecommons.org/licenses/by/4.0/, as recorded in the arXiv licence metadata for this exact version and shown on the abstract page https://arxiv.org/abs/2601.06077v2. Full licence notice: \texttt{licenses/arxiv-2601.06077-CC-BY-4.0.txt}.\par\smallskip

\noindent\textit{Editorial scope.} Counted unit: the article's proposition prop:vo\_coms\_pred (source0.tex lines 631--633). The article's complete original proof of it is delivered with it (source0.tex lines 634--656, one \texttt{proof} environment of 23 lines). No mathematics is written, repaired or re-derived: the statement and the proof are the article's own lines, in the article's own notation, and no retained line is substituted or re-flowed.  Retained uncounted as the source-local context the proof needs: lines 597--604.  Nothing was omitted from any retained span, so the package carries no omission record.  No retained line contains a citation, so the wrapper carries no bibliography and no bibliography entry.  No retained line contains a figure environment, an image-inclusion command or drawing code, so no image is delivered and the package declares no asset dependency.  Editorial formatting only, in addition: an AMS \texttt{amsart} preamble that restates the article's own declarations and macros used by the retained lines, namely the environments \texttt{lemma}, \texttt{proposition} on separate counters, together with the standard packages those lines need.  The retained environments print in this document as Lemma 1, Proposition 1 in order of appearance, whereas the article prints the same items with the numbering of its own full document.  The complete native source of the article as distributed in the arXiv e-print for this exact version is bundled unchanged as \texttt{sources/192384/source0.tex}.
\begin{lemma}\label{lm:D_ell}
Given a loss function $\ell:\sZ\times\sU\rightarrow\R$ and two distributions $P$ and $Q$ on $\sZ$, a decision-theoretic divergence between $P$ and $Q$ with respect to $\ell$ that generalizes the Kullback-Liebler divergence can be defined as
\begin{align}
D_\ell(P,Q) = \E_P[\ell(Z,u_Q)] - \E_P[\ell(Z,u_P)]
\end{align}
where $u_P = \argmin_{u\in\sU}\E_P[\ell(Z,u)]$ and $u_Q = \argmin_{u\in\sU}\E_Q[\ell(Z,u)]$.
This divergence is convex in $P$ when $Q$ is fixed.
\end{lemma}

\begin{proposition}\label{prop:vo_coms_pred}
For a fixed pair of $P_X$ and $P_Y$, and the set of probability transition kernels $\mathcal{M}=\{P_{Y|X}:P_{Y|X}\circ P_X = P_Y\}$ that couple the pair, the value of common sense and the value of prediction both are convex in $P_{Y|X}$ on $\mathcal{M}$.
\end{proposition}

\begin{proof}
First note that the set $\mathcal{M}$ is a linear space, as $(\lambda P_{Y|X} + (1-\lambda)Q_{Y|X})\circ P_{X} = P_Y$ for all $P_{Y|X}, Q_{Y|X}\in \mathcal{M}$ and $\lambda\in[0,1]$.
In addition, the set of joint distributions induced by $P_X$ and $\mathcal{M}$, denoted by $\Pi = \{P_X P_{Y|X}: P_{Y|X}\in\mathcal{M}\}$ is also a linear space, as the marginal distributions of $\lambda P_{X,Y} + (1-\lambda)Q_{X,Y}$ are $P_X$ and $P_Y$ for all $P_{X,Y}, Q_{X,Y}\in {\Pi}$ and $\lambda\in[0,1]$.
It is thus meaningful to discuss convex functions defined on these sets.

The value of common sense can be written as
\begin{align}
\E[\ell(X,Y,\bar u)] - \E[\ell(X,Y,u^*)]  
=   D_{\ell}(P_{X,Y}, P_X P_Y) 
\end{align}
where we have used the definitions of $\bar u$ and $u^*$ and the decision-theoretic divergence in Lemma~\ref{lm:D_ell}.
From the property of the divergence stated in Lemma~\ref{lm:D_ell}, we know that with $P_X$ and $P_Y$ fixed, the value of common sense is convex in $P_{X,Y}$ on $\Pi$; as $P_{X,Y}$ linearly depends on $P_{Y|X}$ when $P_X$ is fixed, the value of common sense is also convex in $P_{Y|X}$ on $\mathcal M$.

The value of prediction can be written as
\begin{align}
& \E[\ell(X,Y,\bar\psi(X))] - \E[\ell(X,Y,\psi^*(X))]  \nonumber \\
= & \sum_{x\in\sX} P_X(x)  \big( \E[\ell(x,Y,\bar\psi(x)|X=x] - \E[\ell(x,Y,\psi^*(x)|X=x] \big) \\
= & \sum_{x\in\sX} P_X(x)  D_{\ell,x}(P_{Y|X=x}, P_Y) 
\end{align}
where in the last step we have used the the fact that $\bar \psi(x) = \argmin_{u\in\sU}  \E[\ell(x,Y,u)]$ and $\psi^*(x) = \argmin_{u\in\sU}  \E[\ell(x,Y,u)|X=x]$ and the definition of the decision-theoretic divergence.
From the property of the divergence stated in Lemma~\ref{lm:D_ell}, we know that with a fixed $P_Y$, $D_{\ell,x}(P_{Y|X=x}, P_Y)$ is convex in $P_{Y|X=x}$ for all $x\in\sX$.
With $P_X$ fixed as well, the value of prediction becomes a nonnegative linear combination of convex functions in $P_{Y|X}$ on $\mathcal{M}$, thus is convex in $P_{Y|X}$ on $\mathcal{M}$.
\end{proof}

\end{document}
